%------------------------------------------------------------------------------%

\usepackage{apresentacao}
\usepackage{mathconfig}
\usepackage{tabto}

\makeatletter
\graphicspath  {{./figs/}}
\def\input@path{{./figs/}}
\makeatother

\newcounter{exemplo}
\setcounter{exemplo}{0}

%------------------------------------------------------------------------------%
\title {Séries de Taylor -- Série Binomial}
\author{\large Luis Alberto D'Afonseca}
\date  {Integrais e Séries}

%------------------------------------------------------------------------------%
\begin{document}

%------------------------------------------------------------------------------%
\begin{frame}
    \titlepage
\end{frame}

%------------------------------------------------------------------------------%
\begin{frame}{Como avaliar funções potência?}
Como avaliar funções potência?
\begin{itemize}[<+->]
  \item \( f(x) = x^r \) \\[3mm]
  \item \( f(x) = \sqrt[3]{x} \) \\[3mm]
  \item \( f(x) = x^{\sfrac{5}{7}} \)
\end{itemize}
\end{frame}

%------------------------------------------------------------------------------%
\section{Série Binomial}

%------------------------------------------------------------------------------%
\begin{frame}{Derivadas da Função Binomial}
Construir a série de Taylor da função \emph{\(\;f(x)=(1+x)^m\;\)} onde $m$ uma é constante
\begin{align*}
  \action<+->{f(x)       & = (1+x)^m                                \\[2mm]}
  \action<+->{f'(x)      & = (1+x)^{m-1}\; m                        \\[2mm]}
  \action<+->{f''(x)     & = (1+x)^{m-2}\; m(m-1)                   \\[2mm]}
  \action<+->{f^{(3)}(x) & = (1+x)^{m-3}\; m(m-1)(m-2)              \\}
  \action<+->{           & \;\;\vdots                               \\
              f^{(k)}(x) & = (1+x)^{m-k}\; m(m-1)(m-2)\cdots(m-k+1)}
\end{align*}
\end{frame}

%------------------------------------------------------------------------------%
\begin{frame}{Derivadas da Função Binomial em Zero}
Avaliando as derivadas em \emph{\(\;x=0\;\)} temos \(\;(1+x)^p=1\;\)
\begin{align*}
  f(0)       & = 1                        \\[2mm]
  f'(0)      & = m                        \\[2mm]
  f''(0)     & = m(m-1)                   \\[2mm]
  f^{(3)}(0) & = m(m-1)(m-2)              \\
             & \;\;\vdots                 \\
  f^{(k)}(0) & = m(m-1)(m-2)\cdots(m-k+1)
\end{align*}
\end{frame}

%------------------------------------------------------------------------------%
\begin{frame}{Coeficientes de Taylor da Função Binomial}
Os coeficientes da série são
\begin{align*}
	c_0 & = 1                        \\[2mm]
	c_1 & = m                        \\[2mm]
	c_2 & = \frac{m(m-1)}{2}         \\[2mm]
	    & \;\;\vdots                 \\
	c_n & = \frac{m(m-1)(m-2)\cdots(m-k+1)}{k!}
\end{align*}
\end{frame}

%------------------------------------------------------------------------------%
\begin{frame}{Coeficiente Binomial}
Para $m$ inteiro positivo

$c_k$ é o número maneiras de escolher $k$ elementos de um total de $m$

\vspace{\baselineskip}
\[
  \frac{m(m-1)(m-2)\cdots(m-k+1)}{k!} \;=\; \frac{m!}{k!(m-k)!} \;=\; \binom{m}{k}
\]
\end{frame}

%------------------------------------------------------------------------------%
\begin{frame}{Coeficiente Binomial}
Por similaridade definimos para $m$ real
\[
  \binom{m}{1} = m
\]
\vspace{0.5\baselineskip}
\[
  \binom{m}{2} = \frac{m(m-1)}{2}
\]
\vspace{0.5\baselineskip}
\[
  \binom{m}{k} = \frac{m(m-1)(m-2)\cdots(m-k+1)}{k!} \qquad k \geq 3
\]
\end{frame}

%------------------------------------------------------------------------------%
\begin{frame}{Série Binomial}
Para \(-1<x<1\) temos
\[ (1+x)^m = 1 + \sum_{k=1}^{\infty} \binom{m}{k} x^k \]
\end{frame}

%------------------------------------------------------------------------------%
\section{Exemplos}

%------------------------------------------------------------------------------%
\addtocounter{exemplo}{1}
\begin{frame}{Exemplo \theexemplo}
  Encontre a Série de Taylor da função $\sqrt[3]{x}$
\end{frame}

%------------------------------------------------------------------------------%
\begin{frame}{Exemplo \theexemplo}
  Escrevemos a função na forma binomial
  \[
    \sqrt[3]{x}
    \pause
    = x^{\sfrac{1}{3}}
    \pause
    = \big(1 + (x-1)\big)^{\sfrac{1}{3}}
    \pause
    = (1+y)^{\sfrac{1}{3}}
  \]
  onde $y=x-1$

  \pause
  \vspace{\baselineskip}
  Usando a Série Binomial
  \[
    (1+y)^{\sfrac{1}{3}} = \sum_{k=0}^\infty \binom{\sfrac{1}{3}}{k}\, y^k
    \qquad y \in(-1, 1)
  \]
\end{frame}

%------------------------------------------------------------------------------%
\begin{frame}{Exemplo \theexemplo}
    Desfazendo a mudança de variáveis
    \begin{align*}
        -1 < y & < 1 \\
        -1 < x-1 & < 1 \\
        0 < x & < 2 
    \end{align*}
    A série converge para $x \in(0, 2)$
\end{frame}

%------------------------------------------------------------------------------%
\begin{frame}{Exemplo \theexemplo}
  \begin{align*}
     \onslide<+->{\sqrt[3]{x}}
     \onslide<+->{& = \sum_{k=0}^\infty \binom{\sfrac{1}{3}}{k} (x-1)^k \\[3mm]}
     \onslide<+->{& = \!\binom{\sfrac{1}{3}}{0} (x-1)^0
       + \!\binom{\sfrac{1}{3}}{1} (x-1)^1
       + \!\binom{\sfrac{1}{3}}{2} (x-1)^2
       + \!\binom{\sfrac{1}{3}}{3} (x-1)^3
       + \cdots \\[3mm]}
     \onslide<+->{& = 1 + \frac{1}{3}(x-1) 
     + \frac{\sfrac{1}{3}\left(\sfrac{1}{3}-1\right)}{2}(x-1)^2 \\[3mm]
     &\hphantom{=1 + \frac{1}{3}(x-1) \;} 
     + \frac{\sfrac{1}{3}\left(\sfrac{1}{3}-1\right)\left(\sfrac{1}{3}-2\right)}{3!}(x-1)^3
       + \cdots \\[3mm]}
     \onslide<+->{& = 1 + \frac{1}{3}(x-1) - \frac{1}{9}(x-1)^2 + \frac{5}{81}(x-1)^3 + \cdots}
  \end{align*}
\end{frame}

%------------------------------------------------------------------------------%
\begin{frame}{\contentsname}
 \tableofcontents
\end{frame}

%------------------------------------------------------------------------------%
\end{document}
%------------------------------------------------------------------------------%